\section{Répétition visible et différenciation généalogique dans une région}
\label{md:ejemplo:region}
Les biographies de \eqref{act21:eq:biografias-tpk-tres-caracteres}
permettent de préciser la différence entre une répétition de chiffres et un
retour de l’état. Pour le canal reconnu comme \(e\), le
calendrier \eqref{act21:eq:calendario-0011-interno} produit
\[
 201101\mid121221\mid102011\mid012222\mid102011.
\]
L’élévation précède la lecture décimale. Les mêmes matrices agissent
dans les deux autres canaux, en conservant leurs émissions initiales et leurs
états régionaux distincts. Cette structure commune permet de comparer
les biographies sans identifier leurs coordonnées.

Soit \(w\) une concaténation de \(L\) trits, et \(N\) sa valeur entière
en base trois. Le cylindre correspondant est
\[
 I_w=[N/3^L,(N+1)/3^L).
\]
Pour \(d\) positions décimales, le préfixe est stable sur tout le
cylindre exactement lorsque
\[
 \left\lfloor\frac{10^dN}{3^L}\right\rfloor
 =\left\lfloor\frac{10^d(N+1)-1}{3^L}\right\rfloor.
\]
C’est l’application du lemme~\ref{act21:censo:estabilidad} au lecteur
décimal. L’extrémité droite semi-ouverte explique le terme \(-1\)
de la division entière. En ajoutant la partie entière de la coordinatisation régionale,
les cinq lectures successives sont
\begin{center}
\begin{tabular}{@{}rl@{}}
\toprule
Longueur tritique & Préfixe stable de la lecture régionale\\
\midrule
6 & \(2.71\ldots\)\\
12 & \(2.71828\ldots\)\\
18 & \(2.7182818\ldots\)\\
24 & \(2.7182818284\ldots\)\\
30 & \(2.71828182845904\ldots\)\\
\bottomrule
\end{tabular}
\end{center}
À la longueur vingt-quatre, on a
\(N=202864003874\) et \(3^{24}=282429536481\).
À la longueur trente, on obtient
\(N=147887858824447\) et \(3^{30}=205891132094649\).
Les divisions entières précédentes certifient les préfixes du tableau.
En particulier, les douze premiers chiffres fractionnaires se regroupent comme
\[
 718281828459
 =7\,\underbrace{1828\,1828}_{\text{réitération positionnelle}}\,459.
\]
Le cylindre de vingt-quatre trits fixe déjà les deux copies et le chiffre
suivant. Cette localisation distingue l’épisode décimal de la
réapparition du bloc tritique \(102011\), situé aux positions
troisième et cinquième : leurs coupures et leurs bases de lecture sont différentes.

La répétition du bloc visible n’identifie pas non plus ses préétats.
Les préétats modulo vingt-sept conservés dans le développement de
cette biographie sont
\[
 (25,11,7,17,8,16),\qquad (7,8,4,23,26,7),
\]
avec pour quotients ultérieurs respectifs
\[
 (6,13,9,2,13,1),\qquad (15,0,3,19,23,25).
\]
L’égalité de l’émission coexiste ainsi avec une généalogie différenciée.
Sa fonction explicative est concrète : l’observation d’un bloc
répété conserve moins d’information que la continuation enrichie
qui le produit. Pour attribuer une différence future exclusivement
à la mémoire, il faut également conserver les phases et les coordinatisations de
comparaison.

Cet épisode relie trois niveaux : les transitions produisent la
chaîne tritique, les cylindres déterminent les chiffres stables et
l’état conserve la provenance d’une émission répétée. La
profondeur arbitraire appartient au prolongement démontré
précédemment ; cette lecture finie précise un épisode de cette
biographie. La régularité qu’il importe d’étudier réside dans la
compatibilité entre les niveaux, non dans le prolongement indéfini
d’un fragment décimal : une troisième copie de \(1828\) commencerait par
\(1\), tandis que la continuation déjà produite commence par \(4\).
